\begin{abstract}
    We introduce \textbf{Sparse Concept Anchoring}, a~method that biases latent space to position a~targeted subset of concepts while allowing others to self-organize, using only minimal supervision (in our setting, labels for $<0.1\%$ of examples per anchored concept). Training combines activation normalization, a~separation regularizer, and anchor or subspace regularizers that attract rare labeled examples to predefined directions or axis-aligned subspaces. The~anchored geometry enables two practical interventions: reversible behavioral steering that projects out a~concept's latent component at inference, and permanent removal via targeted weight ablation of anchored dimensions. Experiments on structured autoencoders show selective attenuation of targeted concepts with negligible impact on orthogonal features, and complete elimination with reconstruction error approaching theoretical bounds. Sparse Concept Anchoring therefore provides a~practical pathway to interpretable, steerable behavior in learned representations.
\end{abstract}

\begin{figure}[h]
    \centering
    \includegraphics[width=0.95\linewidth]{figures/teaser-narrow.png}
    \caption{\textbf{Sparse Concept Anchoring organizes latent space predictably using minimal supervision, enabling behavioral steering and permanent concept deletion.} \textit{Left}: Supervision on \concept{red} during training organizes related concepts around the anchor point. \textit{Center}: The resulting structure enables behavioral steering, demonstrated here by repelling \concept{red} stimuli toward nearby colors while preserving other model capabilities. \textit{Right}: Permanent concept deletion removes \concept{red} responses entirely via weight ablation.}
    \label{fig:teaser}
\end{figure}
